% ============================================================
\chapter{The Simplex Method}
\label{ch:simplex}
% ============================================================

The geometric idea is luminous: since the optimum of a linear program sits at a vertex of the constraint polyhedron, one need only walk from vertex to vertex, improving the objective at each step. This is exactly what the simplex method does, invented by George Dantzig in 1947. Despite a theoretical worst-case exponential complexity (demonstrated by Klee and Minty in 1972), the simplex remains, in practice, one of the most efficient algorithms ever devised---a fascinating theoretical anomaly that still defies our understanding.

\section{Principle of the algorithm}

The fundamental theorem of LP (Theorem~\ref{thm:fundamental_lp}) states that
if an LP has a finite optimum, it is attained at a vertex of the feasible
polyhedron.  The \textbf{simplex method}, due to George Dantzig (1947),
traverses adjacent vertices while improving the objective at each iteration.

\begin{algorithme}[title=\textbf{Simplex overview}]
\begin{enumerate}
  \item Start from an initial vertex (BFS).
  \item Test optimality (reduced costs).
  \item If not optimal, pivot to an improving adjacent vertex.
  \item Repeat until optimality or unboundedness is detected.
\end{enumerate}
\end{algorithme}

\section{Notation and setup}

Consider the LP in standard form:
\[
\begin{aligned}
  \min \quad & z = \mathbf{c}^\top \mathbf{x} \\
  \text{s.t.} \quad & A\mathbf{x} = \mathbf{b}, \quad \mathbf{x} \ge \mathbf{0}
\end{aligned}
\]
with $A \in \R^{m \times n}$, $\text{rank}(A) = m$, $\mathbf{b} \ge \mathbf{0}$.

Partition $A = [B \mid N]$, $\mathbf{x} = (\mathbf{x}_B, \mathbf{x}_N)^\top$,
$\mathbf{c} = (\mathbf{c}_B, \mathbf{c}_N)^\top$:
\[
  B\mathbf{x}_B + N\mathbf{x}_N = \mathbf{b}
  \quad \Rightarrow \quad
  \mathbf{x}_B = B^{-1}\mathbf{b} - B^{-1}N\mathbf{x}_N
\]

\begin{definition}[Reduced costs]
The \textbf{reduced costs} of the non-basic variables are:
\[
  \bar{c}_j = c_j - \mathbf{c}_B^\top B^{-1} \mathbf{a}_j
\]
where $\mathbf{a}_j$ is the $j$-th column of $A$.
\end{definition}

\begin{theorem}[Optimality criterion]
\label{thm:simplex_optimality}
A BFS is optimal if and only if all reduced costs are nonnegative:
$\bar{c}_j \ge 0$ for all non-basic $j$.
\end{theorem}

\section{The simplex tableau}

The simplex tableau organizes all the necessary information:

\begin{center}
\begin{tabular}{c|cccc|c}
  & $x_1$ & $x_2$ & $\cdots$ & $x_n$ & \\
  \toprule
  $z$ & $\bar{c}_1$ & $\bar{c}_2$ & $\cdots$ & $\bar{c}_n$ & $-z$ \\
  \midrule
  $x_{B_1}$ & $\bar{a}_{11}$ & $\bar{a}_{12}$ & $\cdots$ & $\bar{a}_{1n}$ & $\bar{b}_1$ \\
  $x_{B_2}$ & $\bar{a}_{21}$ & $\bar{a}_{22}$ & $\cdots$ & $\bar{a}_{2n}$ & $\bar{b}_2$ \\
  $\vdots$ & $\vdots$ & $\vdots$ & $\ddots$ & $\vdots$ & $\vdots$ \\
  $x_{B_m}$ & $\bar{a}_{m1}$ & $\bar{a}_{m2}$ & $\cdots$ & $\bar{a}_{mn}$ & $\bar{b}_m$ \\
\end{tabular}
\end{center}

where $\bar{A} = B^{-1}A$, $\bar{\mathbf{b}} = B^{-1}\mathbf{b}$.

\section{A simplex iteration}

\begin{algorithme}[title=\textbf{One simplex iteration (minimization)}]
\begin{enumerate}
  \item \textbf{Optimality test}: if $\bar{c}_j \ge 0$ for all $j$,
        \textsc{stop} (optimal solution found).
  \item \textbf{Entering variable}: choose $k = \argmin_j \bar{c}_j$
        (most negative reduced cost).
  \item \textbf{Unboundedness test}: if $\bar{a}_{ik} \le 0$ for all $i$,
        \textsc{stop} (problem is unbounded).
  \item \textbf{Leaving variable}: apply the minimum ratio rule:
        \[
          r = \argmin_{i : \bar{a}_{ik} > 0}
            \frac{\bar{b}_i}{\bar{a}_{ik}}
        \]
  \item \textbf{Pivot}: the pivot element is $\bar{a}_{rk}$.  Perform
        Gauss-Jordan elimination to update the tableau.
\end{enumerate}
\end{algorithme}

\begin{warning}[title=\textbf{Ratio rule}]
Only divide by \textbf{strictly positive} coefficients $\bar{a}_{ik} > 0$.
If all are $\le 0$, the problem is unbounded: $x_k$ can be increased
indefinitely.
\end{warning}

\section{Complete example}

\begin{example}[Simplex step by step]
Solve:
\[
\begin{aligned}
  \max \quad & z = 5x_1 + 4x_2 \\
  \text{s.t.} \quad
    & 6x_1 + 4x_2 \le 24 \\
    & x_1 + 2x_2 \le 6 \\
    & x_1, x_2 \ge 0
\end{aligned}
\]

\textbf{Standard form}: add slack variables $s_1, s_2 \ge 0$ and convert to
minimization ($\min -5x_1 - 4x_2$).

\textbf{Initial tableau}:
\begin{center}
\begin{tabular}{c|cccc|c}
  & $x_1$ & $x_2$ & $s_1$ & $s_2$ & RHS \\
  \toprule
  $z$ & $-5$ & $-4$ & $0$ & $0$ & $0$ \\
  \midrule
  $s_1$ & $6$ & $4$ & $1$ & $0$ & $24$ \\
  $s_2$ & $1$ & $2$ & $0$ & $1$ & $6$ \\
\end{tabular}
\end{center}

\textbf{Iteration 1}: Entering variable $x_1$ ($\bar{c}_1 = -5$).
Ratios: $24/6 = 4$, $6/1 = 6$.  Pivot: row $s_1$, column $x_1$ (pivot = 6).

Divide row 1 by 6, then eliminate $x_1$ from all other rows:
\begin{center}
\begin{tabular}{c|cccc|c}
  & $x_1$ & $x_2$ & $s_1$ & $s_2$ & RHS \\
  \toprule
  $z$ & $0$ & $-\frac{2}{3}$ & $\frac{5}{6}$ & $0$ & $20$ \\
  \midrule
  $x_1$ & $1$ & $\frac{2}{3}$ & $\frac{1}{6}$ & $0$ & $4$ \\
  $s_2$ & $0$ & $\frac{4}{3}$ & $-\frac{1}{6}$ & $1$ & $2$ \\
\end{tabular}
\end{center}

\textbf{Iteration 2}: Entering variable $x_2$ ($\bar{c}_2 = -2/3$).
Ratios: $4/(2/3) = 6$, $2/(4/3) = 3/2$.  Pivot: row $s_2$, column $x_2$
(pivot = $4/3$).

\begin{center}
\begin{tabular}{c|cccc|c}
  & $x_1$ & $x_2$ & $s_1$ & $s_2$ & RHS \\
  \toprule
  $z$ & $0$ & $0$ & $\frac{3}{4}$ & $\frac{1}{2}$ & $21$ \\
  \midrule
  $x_1$ & $1$ & $0$ & $\frac{1}{4}$ & $-\frac{1}{2}$ & $3$ \\
  $x_2$ & $0$ & $1$ & $-\frac{1}{8}$ & $\frac{3}{4}$ & $\frac{3}{2}$ \\
\end{tabular}
\end{center}

All reduced costs are $\ge 0$: optimal!

$x_1^* = 3$, $x_2^* = 3/2$, $z^* = 5(3) + 4(3/2) = 21$.
\end{example}

\section{Initialization: Big-$M$ method}

When the initial tableau does not directly provide a feasible basis ($\ge$
or $=$ constraints), we use \textbf{artificial variables}.

\begin{definition}[Big-$M$ method]
For each equality or $\ge$ constraint (after adding a surplus variable), add
an artificial variable $a_i \ge 0$ to the objective function with a penalty
cost $+M$ (minimization) where $M$ is a large positive number.
\end{definition}

\begin{example}[Big-$M$]
\[
\begin{aligned}
  \min \quad & z = 2x_1 + 3x_2 \\
  \text{s.t.} \quad
    & x_1 + x_2 = 4 \\
    & x_1 + 2x_2 \ge 6 \\
    & x_1, x_2 \ge 0
\end{aligned}
\]
After adding surplus $s_1$ and artificials $a_1, a_2$:
\[
\begin{aligned}
  \min \quad & z = 2x_1 + 3x_2 + Ma_1 + Ma_2 \\
  \text{s.t.} \quad
    & x_1 + x_2 + a_1 = 4 \\
    & x_1 + 2x_2 - s_1 + a_2 = 6 \\
    & x_1, x_2, s_1, a_1, a_2 \ge 0
\end{aligned}
\]

If at optimality $a_1 = a_2 = 0$, the solution is feasible for the original
LP.  Otherwise, the original LP is infeasible.
\end{example}

\section{Initialization: two-phase method}

\begin{algorithme}[title=\textbf{Two-phase method}]
\textbf{Phase I}: Solve the auxiliary problem:
\[
  \min \sum_{i} a_i \quad \text{s.t.} \quad
  A\mathbf{x} + I\mathbf{a} = \mathbf{b}, \quad
  \mathbf{x}, \mathbf{a} \ge 0
\]
\begin{itemize}
  \item If the optimal value is $> 0$: the original LP is \textbf{infeasible}.
  \item If the optimal value is $= 0$: we obtain a BFS for the original LP.
\end{itemize}

\textbf{Phase II}: Use the BFS obtained as a starting point to solve the
original LP with the original objective function.
\end{algorithme}

\begin{bestpractice}[title=\textbf{Big-$M$ vs.\ two-phase}]
\begin{itemize}
  \item The Big-$M$ method is simpler to implement but can cause numerical
        issues if $M$ is poorly chosen.
  \item The two-phase method is numerically more stable and is preferred in
        professional implementations.
\end{itemize}
\end{bestpractice}

\section{Degeneracy and cycling}

\begin{definition}[Degeneracy]
An iteration is \textbf{degenerate} if the minimum ratio is zero: the
leaving variable is already zero and $z$ does not change.
\end{definition}

\begin{theorem}[Cycling risk]
In the presence of degeneracy, the simplex method can theoretically
\textbf{cycle} (loop indefinitely among the same bases).
\end{theorem}

\begin{definition}[Bland's rule]
\textbf{Bland's rule} (smallest-index rule) prevents cycling: choose as the
entering variable the non-basic variable with the smallest index having a
negative reduced cost, and as the leaving variable the basic variable with
the smallest index achieving the minimum ratio.
\end{definition}

\begin{theorem}[Bland's anti-cycling guarantee]
Under Bland's rule, the simplex algorithm terminates in a finite number of
iterations.
\end{theorem}

\section{Simplex complexity}

\begin{proposition}[Number of vertices]
A polyhedron defined by $m$ constraints and $n$ variables has at most
$\binom{n}{m}$ vertices.
\end{proposition}

\begin{remark}
The simplex has $\mathcal{O}(2^n)$ worst-case complexity (Klee-Minty
examples, 1972), but in practice it performs roughly $2m$ to $3m$ iterations
for a problem with $m$ constraints.
\end{remark}

\section{Revised simplex}

The \textbf{revised simplex} avoids maintaining the full tableau.  It
computes only the information needed at each iteration:

\begin{algorithme}[title=\textbf{Revised simplex --- one iteration}]
\begin{enumerate}
  \item Compute $\boldsymbol{\pi}^\top = \mathbf{c}_B^\top B^{-1}$
        (simplex multipliers).
  \item For each $j \notin$ basis, compute
        $\bar{c}_j = c_j - \boldsymbol{\pi}^\top \mathbf{a}_j$.
  \item If $\bar{c}_j \ge 0$ for all $j$: \textsc{stop}.
  \item Choose $k$ with $\bar{c}_k < 0$ (entering variable).
  \item Compute $\mathbf{d} = B^{-1}\mathbf{a}_k$ (direction).
  \item Minimum ratio: $r = \argmin_{i : d_i > 0} \bar{b}_i / d_i$.
  \item Update $B^{-1}$ (product of elementary matrices).
\end{enumerate}
\end{algorithme}

\section{Python implementation}

\begin{codeblock}[title=\textbf{Simplex tableau in Python}]
\begin{minted}{python}
import numpy as np

def simplex_tableau(c, A, b):
    """Solve max c^T x  s.t.  Ax <= b, x >= 0."""
    m, n = A.shape
    # Add slack variables
    tableau = np.zeros((m + 1, n + m + 1))
    tableau[0, :n] = -c          # objective row (min -c^T x)
    tableau[1:, :n] = A
    tableau[1:, n:n+m] = np.eye(m)
    tableau[1:, -1] = b

    basis = list(range(n, n + m))

    while True:
        # Entering variable: most negative reduced cost
        pivot_col = np.argmin(tableau[0, :-1])
        if tableau[0, pivot_col] >= -1e-10:
            break  # optimal

        # Ratios
        col = tableau[1:, pivot_col]
        rhs = tableau[1:, -1]
        ratios = np.full(m, np.inf)
        for i in range(m):
            if col[i] > 1e-10:
                ratios[i] = rhs[i] / col[i]

        pivot_row = np.argmin(ratios) + 1
        if ratios[pivot_row - 1] == np.inf:
            raise ValueError("Problem is unbounded")

        # Gauss-Jordan pivot
        pivot_val = tableau[pivot_row, pivot_col]
        tableau[pivot_row] /= pivot_val
        for i in range(m + 1):
            if i != pivot_row:
                tableau[i] -= (tableau[i, pivot_col]
                               * tableau[pivot_row])
        basis[pivot_row - 1] = pivot_col

    # Extract solution
    x = np.zeros(n)
    for i, var in enumerate(basis):
        if var < n:
            x[var] = tableau[i + 1, -1]

    return x, tableau[0, -1]  # (solution, -z_max)

# Test
c = np.array([5., 4.])
A = np.array([[6., 4.], [1., 2.]])
b = np.array([24., 6.])

x_opt, neg_z = simplex_tableau(c, A, b)
print(f"x* = {x_opt}")
print(f"z* = {-neg_z:.2f}")
\end{minted}
\end{codeblock}

\begin{codeoutput}
x* = [3.  1.5]
z* = 21.00
\end{codeoutput}

\begin{codeblock}[title=\textbf{Verification with SciPy}]
\begin{minted}{python}
from scipy.optimize import linprog

c_min = [-5, -4]
A_ub = [[6, 4], [1, 2]]
b_ub = [24, 6]
res = linprog(c_min, A_ub=A_ub, b_ub=b_ub,
              bounds=[(0, None), (0, None)], method='highs')
print(f"x* = {res.x}, z* = {-res.fun:.2f}")
\end{minted}
\end{codeblock}

\begin{codeoutput}
x* = [3.  1.5], z* = 21.00
\end{codeoutput}

\section{Exercises}

\begin{exercise}[$\star$ --- Simplex tableau]
Solve using the simplex tableau:
\[
\begin{aligned}
  \max \quad & z = 3x_1 + 5x_2 \\
  \text{s.t.} \quad
    & x_1 \le 4 \\
    & 2x_2 \le 12 \\
    & 3x_1 + 5x_2 \le 25 \\
    & x_1, x_2 \ge 0
\end{aligned}
\]
\end{exercise}

\begin{exercise}[$\star\star$ --- Big-$M$]
Solve using the Big-$M$ method:
\[
\begin{aligned}
  \min \quad & z = 4x_1 + x_2 \\
  \text{s.t.} \quad
    & 3x_1 + x_2 = 3 \\
    & 4x_1 + 3x_2 \ge 6 \\
    & x_1 + 2x_2 \le 4 \\
    & x_1, x_2 \ge 0
\end{aligned}
\]
\end{exercise}

\begin{exercise}[$\star\star$ --- Two-phase method]
Re-do the previous exercise using the two-phase method.
\end{exercise}

\begin{exercise}[$\star\star$ --- Degeneracy]
Verify that degeneracy occurs in:
\[
\begin{aligned}
  \max \quad & z = 2x_1 + x_2 \\
  \text{s.t.} \quad
    & x_1 + x_2 \le 4 \\
    & x_1 \le 3 \\
    & x_2 \le 1 \\
    & x_1, x_2 \ge 0
\end{aligned}
\]
\end{exercise}

\begin{exercise}[$\star\star\star$ --- Implementation]
Implement the two-phase method in Python and test it on an LP of your choice
that has $\ge$ and $=$ constraints.
\end{exercise}
